\documentclass[11pt]{article}
\usepackage{amsmath,amsthm,amssymb}
\usepackage[margin=1in]{geometry}
\usepackage[hidelinks]{hyperref}
\newtheorem{theorem}{Theorem}
\newtheorem{lemma}[theorem]{Lemma}
\title{Reciprocal neighbor charges in First-Fit coloring of interval graphs}
\author{kou1353}
\date{November 27, 2026}
\begin{document}
\maketitle

\begin{abstract}
We prove that for every nonempty finite interval graph $G$ of clique number $\omega(G)$, every First-Fit coloring uses at most $(4+2\sqrt{2})\,\omega(G)+\sqrt{2}-2 \approx 6.82843\,\omega(G)-0.58579$ colors, improving the $8\omega(G)-3$ upper bound of Narayanaswamy and Subhash Babu (2008). The proof modifies the ordered column method by granting one unit of income to borrowed rounds, enabling a reciprocal charging inequality across completed neighbor tenures. The rational specialization $35\chi_{\mathrm{FF}}(G)+20\le 239\omega(G)$ (at parameter $a=17/5$) has been formally verified end-to-end in Lean~4.15.0 with no \texttt{sorry} and only standard axioms.
\end{abstract}

\begin{theorem}
For every nonempty finite interval graph of clique number $k$, every
First-Fit order uses a number $m$ of colors satisfying
\[
 m\le (4+2\sqrt2)k+\sqrt2-2.
\]
An entirely rational certificate gives $35m+20\le239k$.
\end{theorem}

\paragraph{Setup.}
Fix the arrival order, the resulting proper coloring, and a closed interval
representation. An interval of color $j$ intersects some interval of each
color $i<j$, since otherwise First-Fit could have assigned $i$.
Choose a finite ordered set of points stabbing all the intervals. All
distinct left endpoints suffice. Write $d$ for the largest number of
intervals containing a chosen point, so $d\le k$. We will prove the theorem
with $d$ in place of $k$. This includes the case in which each interval is
given a point of attachment of depth at most $d$.

\paragraph{The column evolution.}
Choose a real number $a$ with $3<a\le4$. Initially all the chosen points
(columns) are active with balance zero. In round $i$, first determine all
labels simultaneously from the list of columns active before the round.
A column gets income $a$ if it belongs to an interval of color $i$ (real).
Otherwise it gets income $1$ if its immediate active predecessor or
successor is real (borrowed). It gets income zero otherwise (idle).
Rows $i>m$ have no real columns. Charge one unit per round. A real
column survives with balance increased by $a-1$. A borrowed column
survives with the same balance. An idle column of balance greater than
one survives with balance decreased by one. Delete every other idle
column. For accounting on a deletion round, subtract the unit fee there
as well, giving a formal closing balance at most zero. All balances
before deletion rounds are nonnegative.

\begin{lemma}
After round $i\le m$, the surviving points stab every interval of color
at least $i$. The first empty round $N$ is finite and $N\ge m+3$.
\end{lemma}
\begin{proof}
At the start of round $i$, all intervals of color at least $i$ are
covered, by induction. An interval of color $i$ retains every one of its
active points. For a higher interval $I$, take a meeting witness $J$ of
color $i$. Each contains a nonempty consecutive block of the active
points. If the blocks share a point, $I$ retains a real point. Otherwise
the blocks are adjacent: any active point strictly between them would
separate the two intervals. The nearest point of $I$ to the block of $J$
therefore either is real or borrows. This proves coverage. At least one
column is real in round $m$, finishing that round with balance at least
$a-1>2$. It survives the next two idle rounds. After $m$ every remaining
balance decreases until deletion, proving finiteness.
\end{proof}

\paragraph{A terminal column.}
Fix any point $x$ active at the start of round $N$. Let $r,b$ be its
numbers of real and borrowed rounds. Its balance after its $N-1$
surviving rounds is $ar+b-(N-1)\ge0$. Also, its $r$ real rounds have
distinct colors on intervals all containing $x$. Thus
\begin{equation}
 N-1\le ar+b,\qquad r\le d. \label{own}
\end{equation}

\begin{lemma}[Reciprocal charges]
\begin{equation}
 ab+2r\le 2N. \label{recip}
\end{equation}
\end{lemma}
\begin{proof}
Follow the immediate active neighbor of $x$ separately on each side,
always using the order at the start of the round. Partition rounds
$1,\ldots,N$ on each side into maximal consecutive periods with the same
neighbor, or with that side vacant.

An occupied period of length $L$ and occupant $y$ ends on the round in
which $y$ is deleted. Indeed, while $x$ and $y$ both survive, deletion
of other points cannot put a point between them. In the last round every
remaining point is deleted. Let $u\ge0$ be the balance of $y$ before
the period and $S$ its total income in the period. Its formal closing
balance is $u+S-L\le0$, hence $S\le L$. This remains true when $y$
already had a positive balance on becoming a neighbor. Assign the actual
income of $y$ to this side in each of these rounds. On a vacant side
assign one unit just when $x$ is real, and zero otherwise. A vacant
period likewise has assigned sum at most its length. The two side sums
are consequently at most $2N$ in total.

If $x$ borrows, at least one side has a real occupant earning $a$.
If $x$ is real, each occupied side earns at least one: its occupant
either is itself real or sees $x$ as its immediate neighbor and borrows.
A vacant side also contributes one in this case. Real and borrowed
rounds of $x$ are disjoint. Summing these demands gives \eqref{recip}.
\end{proof}

\begin{proof}[Conclusion of the theorem]
Multiply \eqref{own}'s first inequality by $a$ and eliminate $ab$
with \eqref{recip}. This yields $(a-2)N\le(a^2-2)r+a$.
Using $N\ge m+3$ and $r\le d$ gives
\begin{equation}
(a-2)m\le(a^2-2)d+6-2a. \label{family}
\end{equation}
At $a=17/5$, clearing denominators gives $35m+20\le239d$.
The coefficient of $d$ is $4+t+2/t$ for $t=a-2$. It is minimized
at $t=\sqrt2$, by $(t+2/t)^2-8=(t-2/t)^2$. The constant term at
this choice is $\sqrt2-2$. Every order satisfies the result, so
maximizing over orders proves the statement for $\chi_{\mathrm{FF}}$.
The empty graph separately has zero colors and clique number zero.
\end{proof}

\paragraph{Endpoint conventions.}
For a finite family, rank distinct coordinates in increasing order.
At rank $j$ put an open right endpoint at $3j-1$, either closed endpoint
at $3j$, and an open left endpoint at $3j+1$. Convert all resulting
intervals to closed intervals. At a tie two intervals touch precisely
when both endpoints include the point. Strictly ordered endpoints keep
their order. This preserves the intersection graph and every First-Fit
color. Unbounded ends can first be truncated beyond all finite endpoints.
Singleton intervals are allowed. Empty intervals, if admitted, give
isolated vertices. The column evolution keeps the original vertex
colors fixed throughout.

\paragraph{Comparison with the classical count.}
Pemmaraju, Raman and Varadarajan introduced the column method.
Narayanaswamy and Subhash Babu run it on a minimal cover by maximal
cliques, ordered along the line. They keep the first row, labeling holes
there with a dollar sign. Later rows retain occupied columns, holes
beside occupied active columns, and fill cells passing a real-density
test since an adjacency began. The terminal column's dollar signs are
charged to successive neighbors on the two sides. Their disjoint
density-window argument bounds the fill cells. In the all-suffix version
of Joret and Milans, write $\epsilon$ for the density threshold. Up
to end-round constants these estimates read $b\le 2\epsilon N$ and
$r\ge\epsilon(N-b)$. Their product $\epsilon(1-2\epsilon)$ peaks at
$1/8$. The two sides can pay for the same borrowed row, real neighbors
can coincide with a real terminal row, and borrowed rows can lie
inside the windows paying for fill. Each occurrence loses information
in that separate accounting. Tuning the scalar threshold leaves the
coefficient at eight. In the modified evolution a borrowed round
receives a full unit before paying its fee. Symmetry of immediate
neighbors then supplies the $2r$ term in \eqref{recip}. The tenure
estimate depends on the whole period ending at death. Even two disjoint
intervals can give income $a>1$ in the first round of a longer tenure.

The losses in the classical all-suffix count can also be checked as an
identity. At price four, pay zero income on a borrowed round, retaining
the column and flooring its balance at zero on that round. Keep the
same strict survival test on idle rounds. If $P$ borrowed rounds of a
terminal column spend a unit of positive balance, its balance after
round $N-1$ is $3r-f-P=1$, where $f$ is its fill count. Indeed the last
surviving round is an idle fill after the top color, and its positive
integer closing balance is exhausted on round $N$. Since
$r+b+f=N-1$, this gives exactly
\[
8d-m=8(d-r)+(N-m)+(N-2b)+2P.
\]
If $S_L,S_R$ count all real neighbors on each side during their
tenures, then $N-2b$ further splits into
$(N-4S_L)/2+(N-4S_R)/2+2(S_L+S_R-b)$.
The first two terms include unused time and slack at neighbor deaths;
the last term counts redundant neighbor hits. All four terms in the
identity are nonnegative. This explains which information is gained by
changing the income rule.

\paragraph{Linear program.}
At unit normalized lifetime, the homogeneous aggregate constraints are
$ar+b\ge1$ and $2r+ab\le2$. For $D=a^2-2$, the minimum of $r$ in this
relaxation is $(a-2)/D$. A feasible equality point is
$(r,b)=((a-2)/D,(2a-2)/D)$. Dual multipliers $a/D$ and $1/D$ cancel
$b$ and give exactly the asserted lower bound on $r$. This checks
optimality within the two-constraint relaxation, without assuming that
an interval graph realizes the equality point.

\paragraph{Finite-depth consequence.}
Let $A(s)$ be the maximum First-Fit height when each interval has an
attachment point of depth at most $s$, and let $F(k)$ denote the maximum
height at clique number at most $k$. Equation \eqref{family} applies to
the attachment points themselves. The accompanying earlier certificates
give $F(5)\le15$. For $k\ge7$, set $s=k-6$. If the height is at least
$A(s)+1$, consider its top $A(s)+1$ color classes. Some high interval
$I$ has high depth at least $s+1$ everywhere, by the definition of
$A(s)$ and the persistence of lower-color witnesses within that top band.
At most ten lower colors meet the two endpoints of $I$. Delete those
colors, and keep the remaining lower intervals strictly inside $I$.
Every other lower color occurs there because it must meet $I$.
Each retained interval retains all required witnesses of the remaining
lower colors: a witness reaching in from outside would meet an endpoint.
The resulting wall has width at most five. Thus
$F(k)\le A(k-6)+1+10+15$, also valid if the top band did not exist.
The rational bound for $A(k-6)$ gives $35F(k)+544\le239k$ for $k\ge7$.
This optional refinement uses the finite enumeration. The main theorem
is independent of it.

\paragraph{Formal verification.}
The rational theorem $35m+20\le 239d$ has been formalized end-to-end in Lean~4.15.0 (\texttt{ReciprocalColumns.lean}, \texttt{first\_fit\_interval\_rational\_bound}) using only Lean's standard library, with zero \texttt{sorry} and only the standard axioms \texttt{propext}, \texttt{Classical.choice}, and \texttt{Quot.sound}. Accompanying Python scripts (\texttt{audit\_column\_epochs.py}, \texttt{column\_bounds.py}, and \texttt{anchored\_walls.py}) provide independent exact-rational checks. The proof, formalization, and verification scripts are publicly available at \url{https://github.com/kou1353/first-fit-interval-bound}.

\begin{thebibliography}{9}
\bibitem{prv} S.~V. Pemmaraju, R. Raman and K.~R. Varadarajan,
\emph{Buffer minimization using max-coloring}, SODA 2004, 562--571.
\bibitem{nsb} N.~S. Narayanaswamy and R. Subhash Babu,
\href{https://doi.org/10.1007/s11083-008-9076-6}{A note on First-Fit
coloring of interval graphs}, Order 25 (2008), 49--53.
\bibitem{jm} G. Joret and K.~G. Milans,
\href{https://arxiv.org/abs/1006.5704}{First-Fit is linear on posets excluding
two long incomparable chains}, Order 28 (2011), 455--464.
\bibitem{kst} H.~A. Kierstead, D.~A. Smith and W.~T. Trotter,
\href{https://arxiv.org/abs/1506.00192}{First-fit coloring on interval
graphs has performance ratio at least 5}, European Journal of
Combinatorics 51 (2016), 236--254.
\end{thebibliography}
\end{document}
